\Description

在这里描述隐藏的信息、选手需要完成的任务以及相关约束。

\interactor

这是一道交互题。

\InteractionStart

在这里说明交互器首先输出的信息及其范围。

\InteractionQuery

在这里说明每次询问的格式、参数范围和询问次数限制。例如：
\begin{verbatim}
? k
\end{verbatim}

在这里说明交互器返回的信息、数值范围及读取方式。

\InteractionAnswer

在这里说明最终答案的格式。例如：
\begin{verbatim}
! a0 a1 ... an
\end{verbatim}

输出最终答案并刷新标准输出后，程序应立即结束。

\InteractionNotes

每次输出询问后，都必须立即刷新标准输出缓冲区。
例如在 C++ 中可以使用：
\begin{verbatim}
cout << "? " << k << endl;
\end{verbatim}

或：
\begin{verbatim}
cout << "? " << k << '\n' << flush;
\end{verbatim}

不要读取或写入交互协议未规定的信息。
在这里说明超过询问次数、询问格式错误、答案错误的判定，以及协议规定的异常终止方式（如有）。

\InteractionExample

在这里按时间顺序列出完整交互过程，并说明哪些行由交互器发送、哪些行由选手程序发送。
交互记录只用于解释协议，不应作为普通输入与答案样例使用。
